%% LyX 2.5.1 created this file.  For more info, see https://www.lyx.org/.
%% Do not edit unless you really know what you are doing.
\documentclass[brazil]{article}
\usepackage[T1]{fontenc}
\usepackage[utf8]{luainputenc}
\usepackage{babel}
\begin{document}

\section{Economia Marxista 1}

Na parte anterior eu apresentei as bases do meu pensamento. Inclusive de forma sucinte o aspecto econômico da teoria marxista. Agora eu gostaria de aprofundar mais a questão de alguns aspectos da economia marxista com uma abordagem mais técnica. Para isso vou trazer de forma condensada o conteúdo do livro de Paul Singer\cite{singer2000curso} para abordar algumas questões. 

\subsection{Teoria do Valor}

Precisamos começar com o fato de que existe um conflito básico que divide a economia em duas abordagens opostas, esta é uma divisão tão significativa que sequer estas abordagens conservam uma linguagem em comum. A discussão aqui será feita baseado na abordagem que adota a teoria do valor-trabalho. Para entendermos o que é esta teoria precisamos começar discutindo o problema do valor. A economia tem como uma característica fundamental a possibilidade de quantificação, isto é, a maior parte das leis econômicas pode ser expressa matematicamente e verificada empiricamente. Essa possibilidade de quantificação decorre precisamente da teoria do valor.

Este é um conceito básico na economia que permite a utilização de uma unidade essencial, para praticamente todos os fenômenos do mundo econômico. Então o conteúdo desta medida é uma pedra fundamental para todo edifício científico. Existem duas maneiras completamente diferentes de se definir valor: a teoria do valor-utilidade e as diversas teorias do valor-trabalho. Este não é um espaço adequado para compararmos ambas as teorias e por hora, basta que saibamos dessa distinção e que toda discussão feita daqui em diante, quando não dito explicitamente o contrário, se dará fazendo uso da teoria marxista do valor.

Essa teoria expressa o valor a partir das relações sociais de produção estabelecidas entre os homens. O valor neste caso é o fruto das relações que se criam entre os homens na atividade econômica. De maneira mais precisa, se mede pelo tempo de trabalho produtivo que os homens gastam na atividade econômica.

A teoria do valor-trabalho parte da ideia de que a atividade econômica é essencialmente coletiva. Ou seja, não nos interessa as atividades individuais. Se um eletrodoméstico quebra e o próprio dono o consertar, isto é uma atividade particular e do ponto de vista desta teoria, não é uma atividade econômica. Isto é diferente se o dono chamar um eletricista que realizará um serviço de conserto, um serviço que deverá ser remunerado. Podemos ver também que as relações sociais que interessam a esta teoria, no capitalismo são aquelas relações que são mediadas pelo dinheiro.

Segundo esta teoria, o valor de uma mercadoria pode ser medido pelo tempo de trabalho médio socialmente necessário investido nessa mercadoria, desta forma, a teoria do valor-trabalho parte da ideia de que o valor é intrinsecamente social e objetivo, isto é, ele pode ser medido objetivamente. Apesar de que não nos aprofundaremos neste debate é importante ter em mente o caráter social do valor, isto é, o valor de uma dada mercadoria não é dada pelo tempo que esta mercadoria de forma individual levou para ser produzida, mas pela média do tempo que a sociedade em seu conjunto leva para produzir este tipo de mercadoria. Se uma determinada empresa levar o dobro de tempo do que a média para produzir este tipo de mercadoria, sua mercadoria não terá o dobro de valor. Em compensação se outra empresa consegue produzir a mesma mercadoria com metade do tempo médio da sociedade, ela conseguirá produzir mais mercadorias que as empresas concorrentes em um mesmo tempo, de forma que uma vez que cada mercadoria tenha o mesmo valor independente de qual empresa a produziu, esta empresa estará produzindo mais valor que suas concorrentes devido a maior quantidade de mercadoria total que está sendo produzida.

Outro aspecto que não entraremos em detalhes até aqui, é o fato evidente de que as mercadorias devem satisfazer necessidades humanas para terem algum valor. Ainda que esta necessidade não seja suficiente para definir qual é o valor, é necessária para que a mercadoria tenha algum valor. Para a teoria do valor-trabalho, o valor do produto social é medido pelo tempo médio de trabalho socialmente necessário gasto na produção de certa quantidade de mercadorias. Porém este produto social vai encontrar um limite, pois para uma dada população, podemos pensar que ela vai requerer uma determinada quantidade de mercadorias, que vai depender de aspectos da população como cultura, composição etária, do poder aquisitivo, etc... o fato é que qualquer mercadoria produzida além desse limite não é necessária e o trabalho gasto nesta produção não é socialmente necessário, portanto não tem valor.

Evidentemente isto levanta o debate sobre a relação entre o preço e o valor, uma contra-argumentação comum seria de que a quantidade de mercadoria " necessária"{} depende do preço da mesma. A resposta que se encontra na literatura dos partidários do valor-trabalho é de que as mercadorias não chegam ao mercado sem preço. No processo de produção da mercadoria, estas empresas jamais se planejariam para vender suas mercadoria abaixo de um preço mínimo necessário para cobrir os custos e garantir uma margem de lucro, para chegarmos neste preço final no mercado, há uma intrincada série de transições a partir do valor da mercadoria de tal modo que, na abordagem marxista, preço e valor são distintos. Temos então dois sistemas distintos: um para tratar dos preços (sistema de preços) e outro para tratar dos valores (sistema de valores). Não cabe a nós aqui nos aprofundarmos nesta discussão, basta sabermos que existe uma relação indireta que estabelece um vínculo entre valor e preço. Como regra geral, estaremos focando nossa discussão sempre em valor.

Evidentemente existem flutuações no preço (a longo prazo, de acordo com a teoria do valor-trabalho, em torno do valor), mas a existência desta relação (ainda que transformada) entre o preço e o valor, e principalmente devido à própria existência de um conceito objetivo de valor, são características que tornam a teoria do valor-trabalho extremamente interessante e compatível com nossa pesquisa. Como contra-exemplo, podemos lembrar da teoria do valor-utilidade que atribui um conceito subjetivo ao valor.

Alguns breves comentários ainda devem ser feitos antes de concluir a sessão: primeiramente o fato de que teoria do valor-trabalho é histórica por definição, ela tem como fundamento explicar o valor do produto social pela divisão social do trabalho, então evidentemente só pode ter valor na medida em que há esta divisão social do trabalho. Em uma era pré-capitalista onde boa parte dos indivíduos trabalham para si próprio e sua atividade é predominantemente não-econômica, não podemos esperar que esta teoria seja válida. E um segundo comentário diz respeito ao fato de que os adeptos da teoria do valor-trabalho reconhecem a aplicação da teoria do valor-utilidade na prática imediata, na explicação dos preços efetivos no mercado. E por fim, a teoria do valor-trabalho parte da produção, pois considera que o valor não surge no mercado, ela surge no trabalho e é realizado na circulação.

\subsection{Repartição da Renda}

Podemos começar a discussão lembrando que o problema da repartição da renda é um dos tópicos mais antigos e clássicos da Economia Política. David Ricardo considerava este o verdadeiro objeto da economia política. De uma forma geral, podemos dizer que as teorias da distribuição tentam explicar de que maneira o produto social é repartido entre as classes fundamentais da sociedade. Isso equivale a investigar de que forma o produto social é distribuído entre certos rendimentos, classicamente os rendimentos são divididos em: salário, lucro, renda da terra e o juro.

Vamos começar com uma definição do que há para ser repartido, ou seja, vamos definir melhor o que significa o produto social, pra isso, precisamos definir o que é trabalho produtivo e o que é trabalho não produtivo. O primeiro diz respeito ao trabalho que é remunerado, e mais especificamente remunerado por uma parte do capital (denominado capital variável, que consiste exatamente na parte do capital que se transforma em salário). Este é o trabalho responsável por produzir valor e necessariamente deve produzir um produto maior que sua própria remuneração, o que para os partidários desta teoria, é uma determinação social e não técnica. Porém existem outros tipos de trabalho que são trocados por rendimentos e não são produtivos. Podemos entender melhor esta distinção entre trabalho produtivo e não produtivo com um exemplo: pensando em uma cozinheira, se ela trabalha na casa de uma família, ela está recebendo uma parte da renda desta família, apenas transformando o trabalho que presta em renda monetária. Isto é diferente se ela prestar o mesmo trabalho mas agora para um restaurante. Neste novo cenário agora ela passa a desempenhar um trabalho produtivo, pois o trabalho que ela produz entrega ao capitalista que a emprega mais do que ela recebe, agora existe um produto que vai ser transformado em lucro. É apenas neste segundo caso que ela contribui pro produto social.

Um novo exemplo pode ajudar a reforçar esta distinção: se um médico tem uma clínica particular sem funcionários, seu trabalho é não produtivo. Novamente temos apenas a conversão de um trabalho em renda. Porém se o médico trabalha para uma hospital, ele se torna assalariado e entra para o rol dos trabalhadores produtivos, sua remuneração em termos de horas de trabalho socialmente necessárias será necessariamente inferior ao número de horas de trabalho socialmente necessário que ele investiu no serviço. Pois o hospital , como um empreendimento capitalista, necessariamente precisa de uma margem de lucro que será obtida através desta diferença. É também exatamente esta diferença que vai compor o excedente social.

Para esta abordagem econômica, o que interessa é exatamente a divisão do produto social entre entre o produto necessário e o excedente, isto é, como se reparte o produto global entre a parcela necessária para a manutenção da capacidade produtiva, física e mental dos trabalhadores e a outra parte que é o excedente social. Esta repartição se dá principalmente pela luta de classes. Isto é, a determinação do nível salarial do trabalhador não é uma questão intrinsecamente econômica ou técnica, e portanto a própria determinação do produto necessário também não é uma questão puramente técnica. Este nível depende em grande parte da capacidade de luta da classe trabalhadora, por exemplo, se os trabalhadores estão organizados ou não em sindicatos. Em resumo, o salário não tem apenas uma determinação econômica, ele é fortemente dependente do equilíbrio de forças dentro de certos limites.

Evidentemente, a possibilidade de redução da remuneração do trabalho tem limites. Citando limites extremos, podemos comentar que há limites fisiológicos pois abaixo de um certo nível a quantidade e a qualidade da força de trabalho vai reduzir, uma vez que os trabalhadores não conseguem mais viver. E também limites políticos que dependem de circunstâncias históricas e pode levar por exemplo a greves e mobilizações populares de maior ou menor intensidade. Há também evidentemente um limite superior, pois o próprio produto necessário tem como limite o produto social, ou seja, o trabalhador nunca pode receber além daquele ponto em que o excedente social fica igual a zero, mas este limite jamais é alcançado. Antes desse limite ser atingido, ocorre uma restrição da acumulação de capital que conduz a economia à uma crise. Nesta condição, provavelmente os salários retomarão a um patamar ``conveniente'' uma vez que crises aumentam o desemprego e diminuem o poder de barganha do assalariado.

Outro elemento a ser discutido diz respeito à taxa incremental de lucro do capitalista. De acordo com esta teoria, a taxa de lucro é resultado da luta concorrencial entre os capitais. Os ramos que estão dando menos lucro serão aqueles que serão abandonados pelo capital, isto é, há uma tendência de diminuição de fluxo de novos capitais, já os ramos que dando maior lucro terão um afluxo de capital, pois terão a preferência dos capitalistas, e isso fará com que depois de algum tempo a capacidade produtiva deste ramo seja aumentada, levando a uma oferta maior de quantidade de mercadorias. Isto tem também por consequência uma tendência de baixa nos preços da mercadoria, o que em última análise também conduz a uma diminuição da taxa de lucro neste ramo. Esta movimentação do capital produz então uma tendência de equalização da taxa de lucro incremental. Este tendência é desafiada pelo dinamismo tecnológico do sistema o que permite o desenvolvimento cada vez maior das forças produtivas.

Outro ponto que merece destaque é que há um pressuposto de racionalidade nesta teoria. Supõe-se que os capitalistas sabem o que estão fazendo, de modo que o conjunto dos capitalistas tende a agir racionalmente. Apesar disso, um certo número sempre erra e desaparece, e então surgem outros capitalistas que tomam o lugar dos eliminados. 

\subsection{Excedente econômico}

Excedente econômico é em linhas gerais, a parte da produção que não é absorvida pelos gastos necessários à produção da mesma. Na abordagem que estamos adotando coincide com o que é denominado mais-valia total produzido na economia durante um certo tempo. Se denominarmos o produto social como P, o capital constante como c e o capital variável como v, e a mais-valia total como mv, temos então:

\[
P=c+v+mv
\]

\textbf{Capital constante} são os elementos produtivos gastos para obter o produto social, isto é matéria prima, material auxiliar, máquinas e instalações, etc. Enquanto máquinas e instalações compõem o que é chamado de capital fixo, isto é, certos elementos que vão se desgastando durante o processo de produção, a matéria prima faz parte do que é chamado de capital circulante, isto é, o que é totalmente consumido na produção. Capital circulante e capital fixo são subconjuntos do capital constante. Além disso, o capital constante é determinado pela técnica de produção utilizada. O \textbf{capital variável} por sua vez é, como já vimos, a quantidade de salários pago pelo trabalho produtivo, e este por sua vez, como também já vimos é determinado principalmente pela luta de classes, isto é, pelo conflito entre a capacidade que os trabalhadores têm de defender o seu padrão de vida, e a capacidade de que os empregadores têm de reduzir ao máximo a remuneração paga aos trabalhadores.

Ou seja, o excedente econômico, que aparece sob a forma de mais-valia em uma economia capitalista é dado por:

\[
mv=P-(c+v)
\]

Novamente, assim como no caso da discussão da teoria do valor, temos uma grandeza objetivamente determinada. Neste caso, temos que o excedente econômico resulta de uma configuração que é ao mesmo tempo social e técnica, da estrutura produtiva da sociedade. Evidentemente a mais-valia tem duas utilizações possíveis: uma utilização produtiva, ou uma utilização improdutiva. No primeiro caso é um excedente potencial que quando utilizado pela sociedade de forma produtiva, então P aumentará também, por exemplo, aumentando c. Porém no segundo caso, se não for usado produtivamente, então não aumentará.

Mas estamos interessados em como o tamanho do excedente é determinado ao longo do tempo. Há duas formas de aumentar o excedente: com a produção de mais-valia absoluta e a produção de mais-valia relativa. O primeiro caso, como o nome sugere, decorre do fato de que se aumentarmos o trabalho humano gasto durante um certo tempo, sem aumentar a remuneração isso acarretará em um acréscimo de mais-valia. Lembrando que o valor do produto é medido em horas de trabalho socialmente necessário, dessa forma todos os termos da nossa equação também são, esta é a unidade essencial da teoria do valor-trabalho.

Vamos trabalhar com um exemplo prático. Para facilitar os cálculos também vamos considerar trabalhadores com tal eficiência que suas horas de trabalho correspondam exatamente às horas de trabalho socialmente necessárias na sua função, isto é, estes trabalhadores levam o tempo médio de toda sociedade para produzir determinada mercadoria. Então se um trabalhador trabalha 8 horas por dia por 250 dias no ano, isso dá 2000 horas de trabalho. Então 1000 pessoas vão produzir 2 milhões de horas de trabalho socialmente necessárias. Vamos supor que deste valor total, 1.2 milhão de horas socialmente necessárias foi convertido em remuneração da força de trabalho e 800 mil horas é o excedente. Se for possível fazer um funcionário passar a trabalhar 2200 horas por ano, sem aumentar a remuneração, ou seja mantendo constante o capital variável v, então a mais-valia passará de 800 mil horas para 1 milhão de horas, pois todo esse acréscimo será convertido em mais-valia. 

Este acréscimo é a mais-valia absoluta e foi a principal forma de aumentar o excedente utilizado no início da industrialização, porém depois desse início, a indústria passou então a adotar um aumento de mais-valia relativa. Mantendo o exemplo, vamos manter a quantidade de 1.2 milhão de horas socialmente necessárias para manter a força de trabalho de 1000 pessoas. Essa quantidade serve para que essas pessoas se alimentem, se vistam, criem seus filhos, etc... ou seja, esta quantidade está incorporada numa série de bens que permite que vivam e mantenham um dado padrão de vida\footnote{Esta padrão de vida mínimo exigido pelos trabalhadoers, considerado o custo mínimo para a reprodução da sua vida social é determinado historicamente e varia conforme tempo e lugar.}. Se a produtividade aumentar, é possível produzir os mesmos bens de uso em menos horas de trabalho, e este é o objetivo e resultado de todo desenvolvimento tecnológico.

Então agora, digamos que agora com o desenvolvimento tecnológico é possível reduzir o capital variável de 1.2 milhão para 1 milhão de horas, pois agora em cada hora de trabalho, se produz mais bens. Dessa forma, sem reduzir o padrão de vida dos trabalhadores (mas também sem aumentar), cai o montante do capital variável necessário para manter o mesmo padrão de vida para a mesma quantidade de trabalhadores, e consequentemente aumenta o excedente social. Com este aumento o excedente econômico passa a valer também 1 milhão de horas, como no caso anterior, mas agora esta é a mais-valia relativa. Vale lembrar que estamos falando em valor como horas socialmente necessárias, não necessariamente em salário nominal. Um exemplo de quando acontece a mais valia relativa de forma individul é quando um capitalista introduz um avanço tecnológico e consegue aumentar produtividade, dessa forma ele se torna temporariamente capaz de produzir mercadorias em um tempo menor que o tempo socialmente necessário, até que os concorrentes adotem a mesma tecnologia e o tempo socialmente necessário então seja reduzido a um novo patamar, este capitalista consegue maximizar seus lucros.

Em resumo, no caso de mais-valia absoluta o volume total de trabalho socialmente necessário desempenhado pelos trabalhadores aumentou, e no caso da mais valia-relativa o que varia é a distribuição do trabalho entre v e mv, ainda que o volume de trabalho tenha se mantido constante, v reduziu, o que se traduz em um acréscimo em mv, por isso o nome relativo. Historicamente , foi a geração de mais valia relativa a forma que o capitalismo encontrou para aumentar o excedente. Para medir a proporção que o trabalho se reparte entre o capital variável e a mais valia (que juntos são chamados de trabalho vivo, em oposição ao capital constante) foi proposta a taxa de exploração dada por $mv/v$. No exemplo anterior tínhamos $mv/v=800/1200=0.67$. Isso significa que para cada hora de trabalho socialmente necessária o trabalhador trabalha “para si'' , ou seja, que retorna em forma de salário, ele trabalhou mais 40 minutos (0.67h) que será apropriado pelo o capitalista. No caso da mais valia relativa, passamos a ter uma taxa de exploração dada por $mv/v=1000/1000=1$ o que implica que agora para cada hora que o trabalhador “trabalha para si'', ele trabalha outra hora inteira para o capitalista.

Uma característica interessante a ser comentada é que foi observado que a mais valia como proporção de todo o produto se manteve constante a longo prazo, embora flutuasse. Evidentemente isso não implica que o excedente absoluto seja constante, pelo contrário, o excedente tende a ter um crescimento constante, porém principalmente o capital constante tende a crescer também, pois para aumentar a produção é necessário uma quantidade de instalações e máquinas cada vez maior. Uma segunda observação que foi feita é que há uma tendência monopolista do capitalismo.

\subsection{Acumulação e Concentração de Capital}

A acumulação de capital é o processo pelo qual uma parte do excedente econômico é convertido em novo capital, ampliando a sua capacidade de produção. Uma parte do produto social toma a forma física dos meios de vida de trabalhadores e outra de máquinas, matérias primas, prédios e demais infraestrutura que vão se somar aos recursos já existentes. Dessa forma é possível aumentar o nível de produção no período seguinte. Nos interessa saber como esse processo de acumulação de capital se dá.

Vamos começar supondo um capitalismo com um mercado concorrencial, isto é, uma situação em que em um dado ramo qualquer nenhum dos participantes, seja produtor ou comprador, tem capacidade de sozinho determinar o preço. Ainda que alguns autores considerem que o caráter competitivo do capitalismo tende a desaparecer em prol de um caráter monopolista, outros autores argumentam que a medida que se estabelece oligopólios ocorre o oposto: uma intensificação da competição. Vamos considerar o capitalismo ainda em uma fase concorrencial para analisar o processo de acumulação. Neste caso, o estímulo de acumular provém da concorrência, o capitalista precisa usar a mais-valia para acumular porque o sistema capitalista o força a isso, aquele que não cresce e não amplia a sua empresa tende a desaparecer.

O limite da acumulação é atingido então quando o conjunto dos desempregados passa a ser incorporado na economia, uma situação de pleno emprego. Quando a acumulação se acelera, há a possibilidade de que um número cada vez maior de empregos vão ser criados, porém se chegarmos em uma situação de pleno emprego, então o poder de barganha dos trabalhadores aumenta muito. A partir do momento que aumentamos o produto necessário aumentando os salários, então consequentemente reduzimos a mais-valia e reduzimos também então o excedente social, como última consequência temos que a acumulação tende a parar.

Porém, no momento que a economia se aproxima do pleno emprego, as inovações tecnológicas que substituem a mão de obra por máquinas passam a ser altamente rentáveis. Neste caso a acumulação passa a ser muito menos no sentido de estender a capacidade produtiva, e passa a ser muito mais no sentido de aprofundar a capacidade produtiva, ou seja, de mudar a tecnologia visando aumentar a produtividade do trabalho. A consequência desta tomada de decisão é o incremento no índice de desemprego. Este acontecimento recebe o nome de ``desemprego tecnológico''. Com isso recupera-se a existência de um conjunto de desempregados. A passagem do período de acumulação ``extensiva'' para o de acumulação ``'intensiva'' é marcada pela crise.

Porém, em que medida os capitalistas podem investir e aumentar a capacidade produtiva? A capacidade produtiva aumenta cada vez que há um investimento. Mas um investimento só se realiza,na medida em que os produtos gerados pela nova capacidade de produzir são vendidos. Se não forem vendidos a mais valia não se realiza, e dessa forma não pode ser acumulada. Portanto existe a necessidade de uma demanda para que a acumulação se realize, isto é, para que se aumente a capacidade produtiva é preciso ter em vista alguém que vá comprar os produtos adicionais que se vai produzir.

Essa demanda é necessariamente externa a um sistema simplificado em que há apenas trabalhadores e capitalistas, pois enquanto os trabalhadores precisam receber menos do que produzem, e por consequência não podem comprar tudo que foi produzido, os capitalistas são exatamente aqueles interessados em acumular. Esta demanda pode então vir do exterior do sistema capitalista ou do interior, neste último caso isto ocorre fundamentalmente na figura do Estado. O Estado, principalmente no capitalismo contemporâneo tanto nos países em desenvolvimento como nos desenvolvidos têm uma possibilidade de atuar diretamente e indiretamente sobre o nível de acumulação, qualquer teoria que esqueça o Estado está analisando um sistema qualquer que não é relevante para a situação presente do capitalismo.

Relacionado à acumulação, também existe outra tendência central e fundamental no capital que é a concentração de capital. A mera existência do capital da empresa individual já implica em uma concentração de meios de produção, sob o comando único de um proprietário ou um grupo de proprietários. Uma empresa capitalista com um pequeno exército e trabalhadores sob seu comando já implica por si só em uma concentração de recursos produtivos.

A acumulação de capital tende a se acelerar o tempo todo na medida que a economia cresce, uma vez que o sentido do progresso é de aumentar a produtividade humana. Porém este aumento de produtividade do trabalho humano é obtido antes de mais nada através de se colocar à disposição do trabalhador um volume cada vez maior de recursos produtivos, é apenas com o auxílio desta maquinaria que o trabalhador consegue produzir cada vez mais. Este fenômeno é chamado de concentração de capital.

A concorrência entre os capitalistas força-os a adotar a melhor técnica disponível, aquela que é capaz de proporcionar a melhor produtividade e normalmente esta também é a que requer mais capital. Os capitais individuais tendem então a crescer mediante a acumulação de capital, ou seja, através da transformação do excedente em um novo capital que permita que eles não apenas produzam mais, mas produzam mais com menos trabalho. Dessa forma temos que a concentração de capital é resultado da própria acumulação do capital. Este processo obviamente tem um limite que é a própria acumulação da sociedade inteira, os capitais individuais só podem crescer na medida que o capital de toda sociedade cresce. Concluindo, o processo de concentração é então por definição, o crescimento por acumulação dos capitais individuais.

Ainda há um outro processo chamado de centralização que merece ser discutido. Este processo diz respeito à dominação de capitalistas por outros capitalistas. Como a produtividade do trabalho depende do volume de capital posto à disposição do trabalhador, os capitais maiores possuem as condições necessárias para expropriar os menores. O processo de centralização atua como uma força de atração dos capitais maiores sobre os menores. O que resta aos menores é se fundir entre si para poder resistir a pressão dos grandes capitais, transformando-se então eles próprios em um grande capital, ou então ser quebrado e absorvido por grandes empresas.

Em resumo, há duas tendências que devem ser mencionadas: a concentração, que nada mais é que o crescimento da empresa média em função da procura de maior produtividade mediante a acumulação de capital, isto é, da transformação de uma parte dos lucros em novo capital. E a centralização, que decorre diretamente da luta concorrencial e das vantagens das maiores empresas. O primeiro tem como limite a acumulação da riqueza de toda sociedade, e o único limite da segunda é todos os meios de produção estarem concentrados na mão de um único proprietário. É notável novamente a tendência intrínseca do capitalismo em direção ao monopólio.

\subsection{Moeda}

Antes de concluirmos a discussão deste livro, uma última discussão deve ser feita. Enquanto a produção econômica necessariamente aumenta com o tempo, o mesmo não é necessariamente verdade sobre a quantidade de moedas.

O capitalismo exige a existência de uma mercadoria que possa a atuar como equivalente geral das demais. Isto decorre do fato de que no capitalismo cada um dos seus componentes, isto é, as propriedades privadas dos meios de produção, tem autonomia para atuarem de forma desarticulada. Dessa forma é função da troca (ou do mercado, entendido como o lugar onde as trocas se realizam) superar a desarticulação da economia capitalista. O mercado assume então o papel de ajustar as atividades produtivas em um momento posterior à produção e isso é feito por um sistema de prêmios e punições. Aqueles que produzem bens escassos em relação à demanda são recompensados com uma parcela maior da mais-valia, já aqueles que produzem em excesso perdem parte ou toda a mais-valia que esperavam obter.

Toda essa função do mercado se realiza através do processo de circulação e troca de mercadoria, mas para que esse processo de troca possa ocorrer de forma eficiente, é necessário que exista uma mercadoria que possa ser utilizada como equivalente geral às demais. Esta mercadoria passa então servir unicamente como equivalente das demais mercadorias e perde seus outros possíveis usos. Praticamente todas as mercadorias já serviram de moeda alguma vez na história para algum povo, mas para a maior parte da economia capitalista o equivalente geral que foi escolhido foi o metal precioso, especificamente ouro e prata. Dessa forma, a existência da moeda é o elemento central em qualquer economia capitalista, ela desempenha a função de informar aos diferentes produtores a viabilidade econômica da sua atividade pregressa.

Devemos nos ocupar então de saber qual é a quantidade de moeda que deve circular na economia. Como uma mesma unidade monetária pode funcionar para mais de uma transação, a quantidade de moeda necessária na economia não precisa ser a soma total de transações em um dado intervalo de tempo, mas pode ser esse número dividido pelo número médio em que cada unidade monetária é utilizada. Esse número médio se chama de V (velocidade média de circulação da moeda). Vamos a um exemplo: Vamos supor que o volume total de transações durante o ano seja 100 unidades, se cada unidade monetária durante o ano for utilizada em média em 10 transações, a quantidade de moedas existentes precisará ser apenas 100 dividido por 10. Dessa forma, em cada momento existe uma quantidade necessária de moeda (Q) que depende do volume total de transações na economia e da velocidade média de circulação da unidade monetária.

Em um cenário em que o ouro é a moeda, quando há escassez, o preço de todas as mercadorias medidas nesta moeda vai descer. Consequentemente a mesma quantidade de mercadorias, ou melhor, o mesmo valor em mercadorias pode circular com menos moeda. De forma mais precisa, no caso do ouro e outras moeda-mercadorias a inflação e deflação não depende apenas da quantidade da moeda em circulação, mas também das mudanças do seu valor\footnote{Medido da mesma forma que as demais mercadorias: pelo tempo de trabalho socialmente necessário para a sua produção.}, mas por hora isso nos basta pois o caso muda quando se trata da moeda-papel. Especialmente a partir do momento que o estado dentro de uma economia nacional pode impor a aceitação de seus papéis, criando então uma moeda inconversível, isto é,o monopólio da emissão por parte do estado lhe permitiu controlar a quantidade de moeda em circulação independentemente do valor de qualquer moeda-mercadoria, como era o caso do ouro. Na fase atual do capitalismo temos uma relação mais complicada, pois ainda que o estado tenho o monopólio da emissão, ele não tem a capacidade de determinar exatamente a quantidade de moeda pois os bancos também podem criar moeda, mas a complexidade desta discussão foge do escopo desta introdução, então vamos nos concentrar nesta fase anterior ao estado atual do capitalismo.

Neste caso, quando o governo emite, o efeito é similar que quando havia maior produção de ouro no caso da moeda-mercadoria. A quantidade pré-determinada é um certo Q, no momento em que o governo introduz um acréscimo, sem que haja uma alteração na produção de mercadorias, o que ocorre é que este vai reduzir o valor do equivalente em relação às demais mercadorias, e a economia se ajusta através do nível geral dos preços. Com isso cria-se um círculo vicioso, pois no momento em que aumenta-se o valor nominal de todas as transações porque os preços sobem, aumenta-se então a quantidade socialmente necessária de moeda em um momento seguinte.

\bibliographystyle{plain}
\bibliography{tese,Artigo1/references,Artigo2/referencias}

\end{document}
